% AI-produced by GPT-6 Astra (OpenAI), ChatGPT Pro mode, 18 November 2026.
% Root second reading and configuration targeted second reading recorded, 18 November 2026.
% No priority or external/human verification is claimed.
\subsection{Norms of actual point maps over the unbounded nilpotent tower}

Fix a prime $p\geq5$, a prolongation of the $p$-adic valuation to
$\overline{\mathbb Q}$, and write
\[
 K_p=\operatorname{Frac}(\mathbb Z_{(p)}^{\rm sh}),\quad
 I=G_{K_p},\quad N=\bigcap_{c\geq1}\gamma_{c+1}(I),\quad
 K=\overline{\mathbb Q}^{N}=K_{p,\rm nil}.
\]
The lower-central terms are closed.  The quotient $I/N$ is the
maximal pronilpotent quotient.  We keep a fixed complex embedding
$\overline{\mathbb Q}\hookrightarrow\mathbb C$ as well; it is
used only after the arithmetic Kummer statements have been proved.
The $p$-adic and complex absolute values are never identified.
For a field $E$ containing $K$ put
$\kappa_E:E^\times\to H^1(G_E,\mathbb Q_p(1))$.

\begin{claim}{B.B8.1}{Proved}{The three arithmetic inputs at every finite extension}
For every finite extension $E/K$:
\begin{enumerate}
\item $\ker\kappa_E=\mu_\infty$.
\item If $b\in E^\times$, $\alpha\in\mathbb Q_p$, and
      $\kappa_E(b)=\alpha\kappa_E(p)$, then $\alpha\in\mathbb Q$.
\item For every finite extension $E'/E$, complex conjugation in the
      chosen complex embedding does not centralize $G_{E'}$.
\end{enumerate}
These assertions re-prove the inputs A1-B.CF13--14 used below; they
are not deductions from a reconstruction theorem.
\end{claim}
\begin{proof}
The inertia group has a normal wild pro-$p$ subgroup and tame quotient
$\prod_{\ell\ne p}\mathbb Z_\ell(1)$.  In a finite nilpotent
quotient its Sylow subgroups commute.  Consequently
\[
 Q=I/N=Q_{(p)}\times T_{p'},\qquad
 T_{p'}=\prod_{\ell\ne p}\mathbb Z_\ell(1).
\]
The entire tame quotient survives because it is abelian.  The
cyclotomic character factors through $Q$, and on $T_{p'}$ its image
is $\mu_{p-1}\subset\mathbb Z_p^\times$.  All roots of unity lie
in $K$, since their action is already abelian.

Suppose $\kappa_E(a)=0$.  Since $G_E$ fixes all roots of unity,
integral Kummer cohomology is a group of homomorphisms to
$\mathbb Z_p(1)$ and is torsion free.  Rational vanishing therefore
means that all compatible $p$-power roots of $a$ belong to $E$.
Choose a finite Galois number field $D/\mathbb Q$ containing $a$
and algebraic generators of $E/K$, and replace $E$ by $E'=KD$.
This field is stable under the local decomposition group.  Let
$M/\mathbb Q_p$ be the completion of $D$ at the chosen place, made
Galois by the choice of $D$.  Let $a_1,\ldots,a_s$ be the local
Galois conjugates of $a$, and put
\[
 F=K_pD(\mu_{p^\infty}),\qquad
 R=F(a_j^{1/p^m}:1\leq j\leq s,\ m\geq1).
\]
These are Galois extensions of $K_p$, with $R\subset E'$.
The group $D_R=\operatorname{Gal}(R/F)$ embeds in
$\mathbb Z_p(1)^s$ by Kummer theory and is torsion free.  The image
$C_R$ of $N$ in $\operatorname{Gal}(R/K_p)$ is finite, since
$E'/K$ is finite.

Choose $t\in T_{p'}$ with cyclotomic value a primitive
$(p-1)$st root $\zeta$, and a lift $\tau\in I$ of $(1,t)$.
It is central modulo $N$, so its commutator with $D_R$ lies in
$C_R\cap D_R=1$.  Put $J=I_M=I\cap G_D$ and
$J_0=J\cap\ker\chi_p$.  The group $J_0$ maps onto $D_R$.
Conjugation of Kummer roots gives
\[
 \kappa_{J_0}(\tau a)=\zeta\kappa_{J_0}(a).                 \tag{B.NORM.1}
\]
Restriction from $H^1(J,\mathbb Q_p(1))$ to
$H^1(J_0,\mathbb Q_p(1))$ is injective.  Indeed $J/J_0$ is
abelian with nontrivial scalar cyclotomic character, and a cocycle
on an abelian group with such a character is a coboundary: comparing
$c(gh)$ and $c(hg)$ gives
\[
 (\chi(g)-1)c(h)=(\chi(h)-1)c(g).
\]
Choosing $g$ with $\chi(g)\ne1$ proves the assertion.
Restriction from $G_M$ to $I_M$ is injective as well, since
$\mathbb Q_p(1)^{I_M}=0$.  Thus (B.NORM.1) holds in the rational
Kummer group of the finite local field $M$.

Applying the valuation gives $v(a)=\zeta v(a)$, hence $v(a)=0$.
Applying the principal-unit logarithm gives
$\log(\tau a)=\zeta\log(a)$.  The Baker--Brumer theorem says
that $p$-adic logarithms of algebraic units that are linearly
independent over $\mathbb Q$ remain independent over
$\overline{\mathbb Q}$; the complete primary proof opened here
is Rousseau (1969), the unnumbered theorem on printed pages
11-01--11-02.  If $\log(a)\ne0$, this theorem first makes the two
logarithms rationally dependent and then makes their ratio
$\zeta$ rational.  This is impossible for $p\geq5$.
Therefore $\log(a)=0$.  The logarithm has torsion kernel on the
unit group of the finite local field, so $a$ is a root of unity.
The converse is immediate.  This proves (1), with every finite
extension included.

For (2), make the same normal enlargement and write $N'=G_{E'}$,
which is normal in $I$ and finite index in $N$.  Conjugating the
assumed equality on $N'$ by $g\in I$ gives
$\kappa_{E'}(g(b)/b)=0$, because the Kummer class of the rational
number $p$ is invariant.  By (1), each ratio is a root of unity.
The finitely many local conjugates lie in a number field, with a
finite group of roots of unity.  One integer $m>0$ therefore gives
$b^m\in K_p$.

The class $\kappa_I(b^m)-m\alpha\kappa_I(p)$ restricts to zero
on $N'$.  But
$H^1(I/N',\mathbb Q_p(1))=0$.  To see this, its finite normal
subgroup $N/N'$ has zero rational first cohomology and fixes the
coefficient module, so inflation--restriction reduces to $Q$.
The central tame element with cyclotomic value $\zeta\ne1$ makes
every cocycle on $Q$ a coboundary by the preceding displayed identity.
Thus the class was zero already on $I$.  The element $b^m\in K_p$
lies in a finite unramified local extension $M_0/\mathbb Q_p$.
The local restriction to inertia is injective as above.  Applying
its valuation, normalized by $v(p)=1$, gives
$m\alpha=v(b^m)\in\mathbb Z$.  This proves (2).

For (3), first $N\ne1$: the splitting field of $X^p-p$ over
$K_p$ has totally ramified Galois group
$C_p\rtimes\mathbb F_p^\times$, with faithful scalar action,
whose nilpotent residual is nonzero.  The radical degree is $p$
by its value, and the cyclotomic tame degree is $p-1$, so this
calculation does give the full indicated finite quotient of $I$.
The group $N$ lies in wild pro-$p$ inertia, since the tame quotient
is abelian.  It has no torsion, as follows directly from Kummer
theory.  A nonidentity finite-order element would have a power of
prime order $q$; let $H$ be the cyclic subgroup that power generates
and $F=\overline{\mathbb Q}^{H}$.  Then $\mu_\infty\subset F$
and $[\overline{\mathbb Q}:F]=q$.  Kummer theory gives a generator
$\beta$ with $\beta^q=a\in F\setminus F^q$.
Choose $\gamma\in\overline{\mathbb Q}$ with $\gamma^q=\beta$.
Taking the degree-$q$ norm gives
\[
 \operatorname{N}(\gamma)^q
   =\operatorname{N}(\beta)=(-1)^{q-1}a.
\]
The sign is a $q$th power in $F$, since all roots of unity belong
to $F$, contradicting $a\notin F^q$.  Thus $N$ is infinite and
every finite-index subgroup $G_{E'}$ is nontrivial.

The field $E'$ has its nontrivial rank-one henselian valuation.  If
complex conjugation $c$ centralized $G_{E'}$, it would preserve
$E'$.  The transported valuation would be another rank-one henselian
valuation.  If inequivalent, the two are independent, and the
F.~K.~Schmidt theorem (Kuhlmann, Theorem~9.44) would make $E'$
separably closed, a contradiction.  Thus $c$ preserves the
valuation, and uniqueness of prolongation puts $c$ in the local
decomposition group.  Its image in the unramified quotient
$\widehat{\mathbb Z}$ is trivial, and then its image in the
torsion-free tame inertia quotient is trivial.  It would be an
element of order two in wild pro-$p$ inertia, impossible for the
odd prime $p$.  This proves (3).
\end{proof}

\begin{claim}{B.B8.2}{Proved}{Rational residues and norm moduli on affine space}
Let $E/K$ be finite.  A class
$\xi\in H^1(C,\mathbb Z_p(1))$ on a smooth curve over $E$ whose
evaluation at every point over every finite extension is an actual
Kummer value has rational geometric boundary residues.

More generally, let $V\subset\mathbb A^n_E$ be the complement
of a hypersurface, and let
$\eta\in H^1(V,\mathbb Z_p(1))$ have that same actual-value
property.  After a finite extension splitting the irreducible
geometric boundary hypersurfaces, with equations $h_1,\ldots,h_s$,
there are $q_j\in\mathbb Q\cap\mathbb Z_p$ and an arithmetic
class $\beta$ such that
\[
 \eta=\beta+\sum_jq_j\kappa(h_j)
       \quad\hbox{in }H^1(V,\mathbb Q_p(1)).               \tag{B.NORM.2}
\]
If $b(t)$ is any actual value of $\eta$ at
$t\in V(\overline{\mathbb Q})$, then in the fixed complex
embedding
\[
 |b(t)|^2=C\prod_j|h_j(t)|^{2q_j},\qquad C>0,\quad
 C\in\overline{\mathbb Q}.                              \tag{B.NORM.3}
\]
The constant is independent of $t$; a finite extension of point
definition does not change it.
\end{claim}
\begin{proof}
Split the curve boundary and fix a point $q$.  Riemann--Roch gives
a rational function $g$ with sole pole at $q$ and with zeros at
every other boundary point.  Delete zero, infinity and the finitely
many branch values from $\mathbb P^1$.  On the inverse image the
map is finite \emph{\'etale}; corestrict $\xi$ to the resulting
punctured rational line.  Its residues at removed branch values are
zero.  Its residue at infinity is exactly $\operatorname{res}_q\xi$.
The norm formula for a discrete valuation multiplies residues by
residue-field degrees, not by ramification indices; after splitting
the points these degrees are one.  This follows first from
$v(Na)=\sum f_w w(a)$ at finite Kummer levels and then in the
inverse limit.  Hence the corestricted class is an arithmetic class
plus $\alpha\kappa(T)$, with
$\alpha=-\operatorname{res}_q\xi$.

Corestriction evaluates as the product of the field norms of actual
point values, so its evaluations are still actual values.  Choose
$t_0$ and $p^m t_0$ in the punctured line, with $m>0$; there are
only finitely many excluded choices.  Subtracting evaluations gives
an actual value $b$ with
$\kappa_E(b)=m\alpha\kappa_E(p)$.  Claim~B.B8.1(2) makes
$\alpha$ rational.  This independently proves the boundary argument
of A1-B.CF16.

For $V$ the geometric unit group modulo constants is freely
generated by the $h_j$, and its Picard group is zero because its
coordinate ring is a localization of a polynomial UFD.  The Kummer
sequence and the inverse limit therefore give (B.NORM.2) initially
with coefficients in $\mathbb Z_p$.  At a smooth geometric point
of a chosen hypersurface away from the others, take a transverse
line defined over a finite extension.  The restricted class has
the actual-value property, and its residue is exactly the selected
coefficient, with intersection multiplicity one.  The curve result
makes every coefficient rational.  Subtracting the indicated Kummer
classes leaves an arithmetic class, either directly by the same
unit calculation or by the kernel part of Hochschild--Serre.

Choose one algebraic basepoint $t_0\in V$ and a positive integer
$m$ clearing all the $q_j$.  Comparing evaluations at $t$ and
$t_0$ and applying Claim~B.B8.1(1) over a common finite extension
gives
\[
 \left(\frac{b(t)}{b(t_0)}\right)^m
       \prod_j\left(\frac{h_j(t_0)}{h_j(t)}\right)^{mq_j}
       \in\mu_\infty.
\]
Its complex absolute value is one.  Taking the positive $m$th root
proves (B.NORM.3), with
$C=|b(t_0)|^2\prod_j|h_j(t_0)|^{-2q_j}$.
This is positive algebraic because all the numbers in the formula
are algebraic and all the exponents are rational.  The formula
uses one fixed $t_0$; it establishes uniformity of the constant,
rather than choosing a separate constant at each point.
\end{proof}

\begin{claim}{B.B8.3}{Proved}{Reconstructing a norm polynomial without assuming continuity}
Let $X/E$ be smooth geometrically connected affine of dimension
$n\geq1$, and let
$\phi:X_{\rm et}\to\mathbb A^1_{E,\rm et}$ be a pinned geometric
morphism over $E_{\rm et}$.  Write
$f:X(\overline{\mathbb Q})\to\overline{\mathbb Q}$ for its actual
point map.  Suppose it is nonconstant.  For any finite generically
\emph{\'etale} morphism $g:X\to\mathbb A^n_E$ of degree $d$,
the polynomials
\[
 P_t(Z)=\prod_{x\in g^{-1}(t)}(Z-f(x))                    \tag{B.NORM.4}
\]
on its geometric \emph{\'etale} fibres agree, on a nonempty
Zariski open set of algebraic $t$, with a monic polynomial
$P_g\in\overline{\mathbb Q}(t_1,\ldots,t_n)[Z]$ of degree $d$.
No continuity or analyticity of $f$ is assumed.
\end{claim}
\begin{proof}
Relative pinning gives an actual $G_E$-equivariant point map,
including $f(x)\in E(x)$ over every finite field of definition.
This is Carlson--Haine--Wolf, Proposition~3.6, whose hypotheses
are a perfect base field, an algebraic extension of it, and a
target topologically of finite type.  Pulling back the open
complement of any $a\in\overline{\mathbb Q}$ shows that
$S_a=f^{-1}(a)$ is a closed subset after a finite scalar extension.
It is proper when $f$ is nonconstant: if the generic point mapped
to a closed point, continuity and specialization would make the
entire point map constant.  Since $g$ is finite, $g(S_a)$ is a
proper closed subset of $\mathbb A^n$.

Fix $a$.  Choose a principal open $V_a\subset\mathbb A^n$
avoiding $g(S_a)$ and the branch locus, and corestrict the pullback
of $\kappa(Z-a)$ along
$g^{-1}(V_a)\to V_a$.  Its evaluation is the actual value
\[
 b_a(t)=\prod_{x\in g^{-1}(t)}(f(x)-a)=(-1)^dP_t(a).
\]
Indeed its factors lie in the residue fields of the fibre points;
the product is their field norm and belongs to the field of $t$.
Claim~B.B8.2 gives
\[
 |P_t(a)|^2=R_a(t),\qquad
 R_a(t)=C_a\prod_j|h_{a,j}(t)|^{2q_{a,j}},                \tag{B.NORM.5}
\]
where $C_a>0$ is algebraic and uniform in $t$, and all $q_{a,j}$
are rational.  The right side is positive, real analytic, and
algebraic in the independent variables $t,\overline t$ off a
finite union of algebraic hypersurfaces.  Its logarithm is
pluriharmonic there.

We now recover the coefficients from finitely many such equations.
For a monic polynomial $P(Z)=\sum_{k=0}^d c_kZ^k$, $c_d=1$,
the squared modulus $|P(r+is)|^2$ is a polynomial in the real
variables $r,s$, of degree at most $2d$ in each.  Its values on
the grid $0\leq r,s\leq2d$ determine this polynomial by
two-variable Lagrange interpolation.  On changing variables to
$Z=r+is$, $W=r-is$, the coefficient of $Z^kW^d$ is exactly
$c_k$.  Thus each coefficient is a fixed algebraic linear
combination of the $(2d+1)^2$ squared moduli on that grid.
A translated and nonzero algebraically scaled grid has the same
property.  This is a coefficient reconstruction, not a continuity
assumption on the original function.

Choose such a grid $A$, and let $\Omega\subset\mathbb C^n$ be
the complement of the finitely many hypersurfaces needed for its
$R_a$'s.  It is connected.  Applying the interpolation functionals
to these real-analytic functions defines a monic polynomial
\[
 P(t,Z)=Z^d+c_{d-1}(t)Z^{d-1}+\cdots+c_0(t)              \tag{B.NORM.6}
\]
whose coefficients are real analytic on $\Omega$ and algebraic
over $\overline{\mathbb Q}(t,\overline t)$.  At every algebraic
$t\in\Omega$ they equal the coefficients of (B.NORM.4).
For every other algebraic $a$, density of algebraic points and
real-analytic continuation now give
$|P(t,a)|^2=R_a(t)$ wherever both sides are defined.
In particular $\log|P(t,a)|^2$ is pluriharmonic away from its
zeros.  The differential identity extends across the auxiliary
hypersurfaces wherever $P(t,a)\ne0$, because the expression in
the derivatives of $P$ is real analytic there.  Hence this
pluriharmonicity can be used at a fixed $t$ for every algebraic
$a$ except the finitely many roots, and then for every complex
$a$ by density.

Here is the analytic implication, including mixed branches.  For
each $j,k$ put
\[
 Q_{jk}(a)=\frac{P_{t_j\overline t_k}(t,a)}{P(t,a)}
          -\frac{P_{t_j}(t,a)P_{\overline t_k}(t,a)}{P(t,a)^2}.
\]
At a fixed $t$ these are rational holomorphic functions of $a$,
and pluriharmonicity reads
$Q_{jk}(a)+\overline{Q_{kj}(a)}=0$.  A holomorphic rational
function equal to an antiholomorphic one on an open set is
constant.  Since the polynomial is monic, all its differentiated
coefficients have lower $a$-degree, and $Q_{jk}(a)\to0$ as
$a\to\infty$.  Consequently every $Q_{jk}$ is zero.

\paragraph{Local algebraic separation of the root types.}
At a generic real-analytic germ, complexify $\bar t$ to independent
variables $w$. The coefficients of $P$ belong to a finite algebraic
extension of $\overline{\mathbf Q}(t,w)$. In a finite splitting
field the derivations $\partial_{t_j}$ and $\partial_{w_k}$ extend
uniquely, since the characteristic is zero. Decompose $P$ into its
distinct roots $\lambda_\nu$ with positive multiplicities $m_\nu$.
The double and simple partial-fraction poles of
\[
 \partial_{t_j}\partial_{w_k}\log P(t,w,a)=0
\]
give respectively
\[
 m_\nu(\partial_{t_j}\lambda_\nu)
        (\partial_{w_k}\lambda_\nu)=0,
 \qquad
 m_\nu\partial_{t_j}\partial_{w_k}\lambda_\nu=0.
\]
If one $t$-derivative of a root is nonzero, every $w$-derivative
is zero; if one $w$-derivative is nonzero, every $t$-derivative
is zero. This is an identity in a field, so no point-dependent
choice of type is involved. A root with both sets zero is constant.

A root killed by all $w$-derivations is algebraic over
$\overline{\mathbf Q}(t)$. Here is the elementary differential
algebra argument. Take its monic minimal polynomial over
$\overline{\mathbf Q}(t,w)$. Derivations commute with the
automorphisms of a separable normal closure, so all its conjugates
are also killed. Its coefficients, the elementary symmetric
functions in those conjugates, belong to the simultaneous constant
field of the $w$-derivations, namely $\overline{\mathbf Q}(t)$.
The same argument applies with $t$ and $w$ exchanged. A root killed
by both sets is algebraic over $\overline{\mathbf Q}$ and hence
belongs to that algebraically closed field.

We have therefore obtained monic algebraic coefficient germs
\[
 P(t,Z)=H_0(t,Z)A_0(\bar t,Z),
\]
where the roots of $H_0$ are holomorphic or constant, and those of
$A_0$ are antiholomorphic and nonconstant. Repeated roots stay with
their full multiplicities. A holomorphic nonconstant algebraic
germ cannot equal an antiholomorphic one: complexifying such an
equality makes all derivatives in both independent variable sets
zero. This observation gives uniqueness of the division into the
two root types, with constants assigned to $H_0$.

\paragraph{Global continuation, including collisions of opposite types.}
Remove from $\Omega$ a complex algebraic hypersurface containing
the poles and branch loci of the finite algebraic coefficient
extensions of $H_0(t,Z)$ and of the conjugate of
$A_0(\bar t,Z)$. Call the resulting connected complex Zariski
open set $\Omega'$. Each coefficient germ can now be continued
along every path in $\Omega'$. Their product continues to $P$,
by uniqueness of continuation of real-analytic germs. This does
not require the roots of $H_0$ to avoid the roots of $A_0$ on the
real diagonal; continuation of the algebraic coefficients of
the two factors is defined independently of those collisions.

After continuation around any loop, let the resulting monic
factors be $H_1(t,Z)$ and $A_1(\bar t,Z)$. They satisfy
\[
 H_1(t,Z)A_1(\bar t,Z)=H_0(t,Z)A_0(\bar t,Z).
\]
Continuation preserves the vanishing of the relevant derivations
and cannot turn a nonconstant algebraic germ into a constant one.
Uniqueness of the two root types, just proved, forces $H_1=H_0$
and $A_1=A_0$. Thus every coefficient of each factor is single
valued on $\Omega'$; the first is holomorphic there, and the
second becomes holomorphic after conjugating the coordinates.

A single-valued holomorphic function on $\Omega'$ algebraic over
$\mathbf C(t)$ belongs to $\mathbf C(t)$. To see this, normalize
the base in its irreducible function-field extension and remove
the remaining branch and pole locus. This gives a connected
finite \'etale cover. Its complex analytic space is connected:
an irreducible complex variety remains analytically connected
after deleting a proper algebraic subset. A globally defined
branch is a section of this covering and has open-and-closed
image, so the connected cover has degree one. The same reasoning
applies to the antiholomorphic coefficients after conjugation.

Finally a rational function in $\mathbf C(t)$ algebraic over
$\overline{\mathbf Q}(t)$ belongs to
$\overline{\mathbf Q}(t)$. One concrete descent argument is to
bound numerator and denominator degrees and choose sufficiently
many algebraic parameter points avoiding their poles and the
degenerate specializations of the algebraic equation. The values
are algebraic. Rational interpolation, normalized by one nonzero
coefficient, then recovers the function by a finite linear system
with algebraic entries. Normalize one denominator coefficient
known to be nonzero. A solution exists over $\mathbf C$, hence
also over $\overline{\mathbf Q}$ by Gaussian elimination.
Sufficiently many interpolation points determine the rational
function uniquely (cross-multiply and interpolate a polynomial
of twice the degree bound), so this algebraic solution represents
the original function. Consequently
\[
 \begin{gathered}
 P(t,Z)=H(t,Z)A(\overline t,Z),\\
 H\in\overline{\mathbf Q}(t)[Z],\qquad
 A\in\overline{\mathbf Q}(\bar t)[Z].
 \end{gathered}\tag{B.NORM.8}
\]

It remains to exclude a nonconstant $A$.  Choose a finite number
field $M$, stable under complex conjugation, containing the
coefficients of $H,A,g,X$ and all auxiliary hypersurfaces, and
work with $\sigma\in G_{EM}$.  Equivariance of the actual point
map gives $P_{\sigma t}=\sigma(P_t)$ at every algebraic parameter
where defined.  Cancelling the common monic factor $H(\sigma t,Z)$
in (B.NORM.8) gives
\[
 A(c\sigma t,Z)=A(\sigma c t,Z),                        \tag{B.NORM.9}
\]
where $c$ denotes coordinatewise complex conjugation.  If a
coefficient $a(w)$ of $A$ is nonconstant, put
$\tau=c\sigma c\sigma^{-1}$.  This automorphism fixes $M$, and
(B.NORM.9) says $a(\tau u)=a(u)$ on a nonempty Zariski open set
of algebraic $u$.  Since $\tau$ fixes the coefficients,
$\tau(a(u))=a(u)$.  A nonconstant rational function on affine
space has cofinite image in the algebraically closed field
$\overline{\mathbb Q}$ (restrict to a suitable line if necessary).
Thus $\tau$ fixes a cofinite subset, and then fixes every element
by subtraction of two elements in that subset.  Hence $\tau=1$.
This would hold for every $\sigma\in G_{EM}$, contradicting
Claim~B.B8.1(3).  All coefficients of $A$ are therefore constant;
by the convention on its roots it has degree zero.  The polynomial
$P=H$ is rational in $t$, as required.
\end{proof}

\begin{claim}{B.B8.4}{Proved}{Generic norms force one algebraic graph}
In the setting of Claim~B.B8.3 there is a unique
$h\in E(X)$ such that $f(x)=h(x)$ at every algebraic point of
some nonempty Zariski open subset of $X$.  The rational function
is regular everywhere and agrees with $f$ at every algebraic
point.
\end{claim}
\begin{proof}
Choose one finite Noether normalization $g$.  Claim~B.B8.3 gives
a monic rational equation $P_g(g(x),f(x))=0$ on a dense open.
After shrinking to a smooth affine $X_0$, all its coefficients are
regular.  The Zariski closure $Z$ of the graph of $f$ over $X_0$
is therefore contained in a scheme finite over $X_0$.
This proves algebraicity of the containing correspondence before
any assertion about the graph itself.  Let $Z_1,\ldots,Z_r$ be
the irreducible components dominating $X_0$, and let their degrees
over $X_0$ be $e_i\geq1$.  There is at least one such component,
since the graph has a point over every algebraic point of $X_0$.
The actual graph is dense in each $Z_i$ after the other components
have been removed: otherwise its closure would not have that
irreducible component.  Since $Z_i$ dominates $X_0$, these graph
points project to a dense subset of $X_0$.  Thus the coordinate
$y$ is nonconstant on each $Z_i$: a component with $y=a$ would
make $S_a=f^{-1}(a)$ dense in $X_0$, whereas $S_a$ is a proper
closed subset for a nonconstant pinned point map.

We choose a second finite normalization
$g':X_0\to\mathbb A^n$ with the following additional property:
each $(g',y):Z_i\to\mathbb A^n\times\mathbb A^1$ is birational
onto its image, and these images are pairwise distinct.  Here is a
generic-projection proof.  Embed $X_0$ in affine space.  Over the
function field $\overline{\mathbb Q}(y)$ the union of the generic
fibres of the $Z_i$ is a reduced union of $(n-1)$-dimensional
subvarieties of that affine space: a fibre of
$Z\subset X_0\times\mathbb A^1$ embeds in $X_0$ when $y$ is
fixed.  A generic linear projection to $\mathbb A^n$ is
generically one-to-one on that union.  The incidence proof counts
pairs of distinct points, of dimension $2(n-1)$, and the $n$
independent conditions that their images coincide; the resulting
double locus has dimension at most $n-2$, less than that of a
component.  The characteristic is zero, so generic degree one is
birational.  The same estimate separates different components.
For $n=1$ the generic fibres are finite reduced sets, and a generic
linear function separates them directly.  These dimension estimates
are on the fibres over the generic value of $y$.  The conditions
first give a nonempty open subset of coefficient space over
$\overline{\mathbb Q}(y)$.  It contains an open defined by
finitely many nonvanishing polynomials in the coefficients whose
coefficients are rational in $y$.  After clearing denominators
and expanding in $y$, failure of one such nonvanishing identically
in $y$ is the simultaneous vanishing of finitely many polynomials
over $\overline{\mathbb Q}$.  Its complement is a nonempty open
in coefficient space over $\overline{\mathbb Q}$.  Thus these
conditions can be imposed simultaneously with the nonempty open
condition of finite Noether normalization of $X_0$.  They have a
choice over a finite extension of $E$, as permitted throughout.

Write $d'=[\overline{\mathbb Q}(X_0):
\overline{\mathbb Q}(g')]$.  Birationality of $(g',y)$ says
\[
 \overline{\mathbb Q}(Z_i)=
 \overline{\mathbb Q}(g',y).
\]
The minimal polynomial of $y$ over
$\overline{\mathbb Q}(g')$ therefore has degree $e_i d'$.
These minimal polynomials are distinct for different components.
On the other hand Claim~B.B8.3 supplies a monic norm polynomial
$P_{g'}$ of degree $d'$ vanishing on the actual graph, hence on
all its dominating closure components.  Their distinct minimal
polynomials divide it.  Thus
\[
                    d'\sum_i e_i\leq d',
\]
so $r=1$ and $e_1=1$.  The graph agrees with one rational function
$h$ off a proper closed subset.  Uniqueness on a dense open and
the $G_E$-equivariance of $f$ imply $\sigma h=h$ for every
$\sigma\in G_E$, hence $h\in E(X)$.

We now return to the original smooth affine $X$ and remove both
possible poles and exceptional point values; agreement on a dense
set alone would not do this.  Enlarge the exceptional subset to a
proper closed $B\subset X$ that contains $X\setminus X_0$ and
outside which agreement and regularity hold.  Fix any algebraic
point $x_0\in B$.  After a finite scalar extension take a finite
linear normalization $g':X\to\mathbb A^n$ of this original $X$
such that
$t_0=g'(x_0)$ is an \emph{\'etale} fibre and
\[
                g'^{-1}(t_0)\cap B=\{x_0\}.
\]
Such projections exist after enlarging the affine embedding by
polynomial functions of degree at most a fixed sufficiently large
integer.  These functions separate values and first jets at any
two distinct points.  Impose $g'(x_0)=0$.  At each other point
$x$, the conditions $g'(x)=0$ have codimension $n$ in the
parameter space of the $n$ functions; requiring also a singular
differential has one further independent determinantal condition,
because first jets can be prescribed.  The incidence over the
$n$-dimensional $X\setminus\{x_0\}$ therefore has image of
positive codimension.  The incidence over $B\setminus\{x_0\}$,
with just the value conditions, also has image of positive
codimension since $\dim B<n$.  At $x_0$ invertibility of the
differential is a nonempty open condition.  The highest-degree
forms can simultaneously be chosen with no common zero on the
projective boundary, which is the nonempty open condition making
the map finite.  Thus a parameter outside these finitely many
proper closed subsets gives exactly the asserted fibre.  This
also supplies a direct jet-and-incidence proof of the required
smooth-fibre Bertini step.

Choose the interpolation grid in Claim~B.B8.3 to avoid the finitely
many actual values of $f$ on this fibre.  For each grid value $a$,
the closed set $g'(S_a)$ does not contain $t_0$, so the principal
open $V_a$ used in that proof can be chosen to contain $t_0$.
The reconstructed coefficients are then real analytic near $t_0$,
are rational by Claim~B.B8.3, and have no pole at $t_0$.
Their polynomial is equal generically to the ordinary norm
characteristic polynomial of $h$, since $f=h$ off $B$.
Consequently $h$ satisfies a monic equation with coefficients
regular at $t_0$.  The local rings of $X$ are normal, so integrality
makes $h$ regular at every point of that fibre, including $x_0$.

At $t_0$ the polynomial equals the product of $(Z-f(x))$ by the
finite-grid reconstruction, and equals the product of $(Z-h(x))$
by specialization of the finite \emph{\'etale} norm.  Every factor
except the one at $x_0$ is already the same.  Cancelling them gives
$f(x_0)=h(x_0)$.  Since $x_0$ was arbitrary, this proves both
global regularity and equality at all geometric points.
\end{proof}

\begin{claim}{B.B8.5}{Proved}{Pinned reconstruction over the nilpotent union for $p\geq5$}
For every $p\geq5$ and every finite extension $E/K_{p,\rm nil}$,
the functor
\[
 \operatorname{Sch}^{\rm ft}_E[\operatorname{UH}^{-1}]
       \longrightarrow\operatorname{RTop}^{\rm pin}_E,
 \qquad X\longmapsto X_{\rm et}
\]
is fully faithful.  Here the target uses geometric morphisms over
$E_{\rm et}$, with their relative structural isomorphisms, that
take closed points to closed points.  Equivalently, for finite-type
$E$-schemes with seminormal source, pinned relative geometric
morphisms of their \emph{\'etale} topoi come from unique morphisms
of $E$-schemes.
\end{claim}
\begin{proof}
Claims~B.B8.3--B.B8.4 construct an actual morphism
$X\to\mathbb A^1$ agreeing with any pinned point map when $X$
is smooth geometrically connected affine.  A constant point map is
the constant morphism; if needed a rational point ensures its value
lies in the base field.  For a smooth connected source that is not
geometrically connected, pass to a finite extension splitting its
geometric components, apply the construction to each, and descend
the unique resulting morphism using equality of its geometric
point values.  Every arithmetic assertion used above is stable
under those finite extensions by Claim~B.B8.1.

We apply the exact reduction theorem of Carlson--Haine--Wolf
(2024), Theorem~4.18, implication (4.18.4)$\Rightarrow$(4.18.1).
It holds for \emph{any} field, without their later amenability
hypothesis.  It requires precisely an actual morphism agreeing on
geometric points for each affine regular connected source with a
rational point and target $\mathbb A^1$.  In characteristic zero
regular finite-type schemes are smooth, and absolute weak
normalization of a regular scheme is itself.  Thus the requirement
has been met.  Their proof uses h-descent and reconstructs ring
operations from geometric-point evaluations; in particular it
identifies the entire inverse-image functor, not merely the point
map.  Finally their Corollary~2.32 identifies the localization at
universal homeomorphisms with the stated seminormal-source
formulation in characteristic zero.  This closes the
arbitrary-source and pointwise-trivial-functor gaps left by
A1-B.CF15--17 for these fields.
\end{proof}

\paragraph{What this changes, and what it does not decide.}
The new step is the use of all finite norm polynomials and a generic
projection to eliminate a multivalued algebraic correspondence.
It handles the proper-cohomology component on an arbitrary source
by corestriction to a principal open in affine space, whose Picard
group is zero.  It does not assume that this component vanished
on the original curve.  The finite interpolation script checks
the coefficient reconstruction for sample degrees only; all
degrees and the profinite arithmetic claims are proved above.

The remaining B8 list is not empty: other mixed-characteristic
henselian fields not covered by the proved positive and negative
classes still require further work.
The maximal metabelian fields and every finite extension are
settled subsequently by Claim~B.META.14; B.138 gives a uniform
all-prime proof for the nilpotent unions, including $p=2,3$.
Those later proofs do not use the torsion-only kernel needed
in the present argument.  Here we recheck the two small-prime
counterexamples to that arithmetic input.
For $p=2$, the splitting field of $X^{2^m}-2$ over $K_2$ has
a $2$-group as Galois group: the cyclotomic and radical degrees
are both powers of two.  It is nilpotent, so every such root
belongs to $K_{2,\rm nil}$.  Thus the nontorsion element $2$
has zero rational $2$-Kummer class there.

For $p=3$, set $z=\sqrt{-3}$ and $a=(2+z)/(2-z)$.
The quadratic extension $F_0=K_3(z)=K_3(\mu_3)$ sends $a$
to $a^{-1}$.  The extension
$R_m=F_0(\mu_{3^m},b)$, $b^{3^m}=a$, is Galois over $K_3$
with a normal $3$-subgroup $P$ and quotient $C_2$.
Cauchy's theorem supplies an involution $s$; it inverts
$\mu_{3^m}$ and has $s(b)=\xi b^{-1}$.  Its square gives
$\xi^2=1$, hence $\xi=1$ since $\xi\in\mu_{3^m}$.
If $g\in P$ has $g(b)=\eta b$ and
$g(\zeta_{3^m})=\zeta_{3^m}^{u}$, direct substitution gives
$sgs^{-1}(b)=\eta b$ and
$sgs^{-1}(\zeta_{3^m})=\zeta_{3^m}^{u}$.
It follows that $s$ centralizes $P$ and
$\operatorname{Gal}(R_m/K_3)=P\times C_2$ is nilpotent.
Thus every $3$-power root of $a$ belongs to $K_{3,\rm nil}$.
The equality $a+a^{-1}=2/7$ proves that $a$ is not a root
of unity, since the sum of a root of unity and its inverse is
an algebraic integer.  These calculations re-prove the two
exceptions in A1-B.CF13.  They refute input~(1) at the two small
primes, not reconstruction itself.  No such field is promoted by
removing $p\geq5$ from the proof.

The current Carlson--Stix revision, arXiv:2410.22474v2 of
20 September 2026, was also opened.  Its Theorem~A/Theorem~4.1
assumes a sub-$p$-adic field.  Its new relative
Theorem~D/Theorem~6.12 assumes a geometrically unibranch
noetherian Nagata base whose component function fields are
sub-$p$-adic.  Neither directly supplies the conclusion above.
Indeed $K_{p,\rm nil}$ contains $\mu_\infty$, whereas the
algebraic constant field of a finitely generated extension of a
$p$-adic local field has only finitely many roots of unity.
Its point-theoretic open-or-constant Proposition~3.17 requires
Kummer faithfulness, which fails here on $\mu_\infty$.

\paragraph{Checks and provenance.}
These are new proofs in this continuation, written by GPT-6 Astra
(OpenAI), ChatGPT Pro mode, on 18 November 2026.  They make no
claim of priority or human verification.  The Kummer inputs were
re-proved from the local inertia and logarithm arguments; the
norm/interpolation and generic-projection steps are separate from
the earlier affine-line proof.  A second reading, and its exact
scope, must be recorded before these working statements are
treated as checked in the final register.
